表达式 \co@subspan(M, K).at(i, j)@ 使用分块坐标。没有显式步长时，元素起点为 \co@(i * M, j * K)@。指定步长会改变相邻分块起点之间的距离，但不改变选中分块的大小。

加上 \co@.step(sM, sK)@ 会改写步长：分块之间相隔 \co@sM@ 行、\co@sK@ 列，可以按指定间隔排列：

\begin{lstlisting}[language=croqtile,caption={带步长的分块：.subspan、.step 与 .at}]
matrix.subspan(16, 16).step(32, 32).at(i, j);
\end{lstlisting}

图~\ref{fig:f09-fig5_subspan_step} 对比了紧凑排布与带步长的分块。

\begin{figure}[htbp]\centering
\includegraphics[width=\linewidth,height=0.56\textheight,keepaspectratio]{f09-fig5_subspan_step_light.png}
\caption{小图比较紧密排列的 [3, 3] 分块与 step(5, 5) 的间隔排列。正文的 [16, 16] 分块采用 step(32, 32)，留下 16 个元素的间隔。}
\label{fig:f09-fig5_subspan_step}
\end{figure}

分块 \co@(0,0)@ 从 \co@(0,0)@ 开始，而分块 \co@(1,0)@ 从第 32 行开始（不是第 16 行），留出 16 行的间隔。省略 \co@.step@ 时，步长等于分块尺寸，也就是紧凑排列的情形。

步长小于分块尺寸时，窗口相互重叠；步长更大时，窗口之间留下空隙。例如，16 行窗口每次前进八行，会与下一个窗口共享一半行。步长控制起点序列，不应与分块内部元素之间的存储步幅混淆。
